% =========================================================
% Models, Theories, Logical Consequence, Logical Equivalence
% =========================================================

\subsection{Semantics: Models, Theories, and Logical Consequence}

\begin{remark*}[Section toolkit: Minimal model notation]
Volume I uses theories only as sets of sentences so that the notation
\(\mathcal M\models T\) and \(T\models\varphi\) can be read. The systematic
study of theories and their classes of models belongs to Volume VIII Model
Theory.
\end{remark*}

\begin{definitionbox}{Definition (Model of a Sentence)}
\begin{definition}[Model of a Sentence]
\label{def:model}
Let \(\varphi\) be a sentence and let \(\mathcal M\) be a structure for the
language of \(\varphi\). The structure \(\mathcal M\) is a \emph{model of}
\(\varphi\) exactly when \(\mathcal M\models\varphi\).
\end{definition}
\end{definitionbox}

\begin{remark*}[Standard quantified statement]
\[
\mathcal M \text{ is a model of } \varphi
\Longleftrightarrow
\mathcal M\models\varphi.
\]
\end{remark*}

\begin{remark*}[Predicate reading]
\(\operatorname{ModelOfSentence}(M,\varphi)\) means that \(M\) is a model of
the sentence \(\varphi\).
\end{remark*}

\begin{remark*}[Negated quantified statement]
\[
\mathcal M \text{ is not a model of } \varphi
\Longleftrightarrow
\mathcal M\not\models\varphi.
\]
\end{remark*}

\begin{remark*}[Negation predicate reading]
\(\neg\operatorname{ModelOfSentence}(M,\varphi)\) means that \(\varphi\) is not
true in \(M\).
\end{remark*}

\begin{remark*}[Failure modes]
\begin{description}
  \item[Exposition.] The only failure mode is semantic: the sentence is false in the proposed
structure.
\end{description}
\end{remark*}

\begin{remark*}[Failure modes]
\begin{description}
  \item[Exposition.] \[
\underbrace{\mathcal M}_{\text{proposed structure}}
\qquad
\underbrace{\mathcal M\not\models\varphi}_{\text{sentence false in the structure}}.
\]
\end{description}
\end{remark*}

\begin{remark*}[Interpretation]
A model is not an extra kind of object. It is a structure viewed relative to a
sentence that it makes true.
\end{remark*}

\begin{dependencies}
\begin{itemize}
  \item \hyperref[def:truth-structure]{Truth in a Structure}
\end{itemize}
\end{dependencies}

\begin{definitionbox}{Definition (Model of a Set of Sentences)}
\begin{definition}[Model of a Set of Sentences]
\label{def:model-theory-set}
Let \(T\) be a set of sentences and let \(\mathcal M\) be a structure for their
common language. The structure \(\mathcal M\) is a \emph{model of} \(T\),
written \(\mathcal M\models T\), exactly when
\(\mathcal M\models\theta\) for every \(\theta\in T\).
\end{definition}
\end{definitionbox}

\begin{remark*}[Standard quantified statement]
\[
\mathcal M\models T
\Longleftrightarrow
\forall \theta\in T\;(\mathcal M\models\theta).
\]
\end{remark*}

\begin{remark*}[Predicate reading]
\(\operatorname{ModelOfTheory}(M,T)\) means that \(M\) is a model of every
sentence in \(T\).
\end{remark*}

\begin{remark*}[Negated quantified statement]
\[
\mathcal M\not\models T
\Longleftrightarrow
\exists \theta\in T\;(\mathcal M\not\models\theta).
\]
\end{remark*}

\begin{remark*}[Negation predicate reading]
\(\neg\operatorname{ModelOfTheory}(M,T)\) means that at least one sentence of
\(T\) is false in \(M\).
\end{remark*}

\begin{remark*}[Failure modes]
\begin{description}
  \item[Exposition.] A proposed model of a set of sentences fails by a single sentence of the set
that is false in the structure.
\end{description}
\end{remark*}

\begin{remark*}[Failure modes]
\begin{description}
  \item[Exposition.] \[
\underbrace{\exists \theta\in T}_{\text{failed sentence}}
\qquad
\underbrace{\mathcal M\not\models\theta}_{\text{not true in the structure}}.
\]
\end{description}
\end{remark*}

\begin{remark*}[Interpretation]
The notation \(\mathcal M\models T\) abbreviates a universal check over the
sentences in \(T\). No additional theory of model classes is developed here.
\end{remark*}

\begin{dependencies}
\begin{itemize}
  \item \hyperref[def:model]{Model of a Sentence}
\end{itemize}
\end{dependencies}

\begin{definitionbox}{Definition (Theory)}
\begin{definition}[Theory]
\label{def:theory}
In Volume I, a \emph{theory} is a set of sentences in a fixed first-order
language.
\end{definition}
\end{definitionbox}

\begin{remark*}[Standard quantified statement]
\[
T \text{ is a theory}
\Longleftrightarrow
\forall \theta\in T\;(\theta \text{ is a sentence of the fixed language}).
\]
\end{remark*}

\begin{remark*}[Predicate reading]
\(\operatorname{Theory}(T)\) means that \(T\) is a set of sentences in a fixed
first-order language.
\end{remark*}

\begin{remark*}[Negated quantified statement]
\[
\exists \theta\in T\;(\theta \text{ is not a sentence of the fixed language}).
\]
\end{remark*}

\begin{remark*}[Negation predicate reading]
\(\neg\operatorname{Theory}(T)\) means that \(T\) is not a set consisting only
of sentences of the fixed language.
\end{remark*}

\begin{remark*}[Failure modes]
\begin{description}
  \item[Exposition.] A proposed theory fails in this minimal sense when it contains something other
than a sentence of the fixed language.
\end{description}
\end{remark*}

\begin{remark*}[Failure modes]
\begin{description}
  \item[Exposition.] \[
\underbrace{\theta\in T}_{\text{member of the proposed theory}}
\qquad
\underbrace{\theta \text{ is not a sentence}}_{\text{bad member}}.
\]
\end{description}
\end{remark*}

\begin{remark*}[Interpretation]
This is intentionally modest. Complete theories, elementary classes,
categoricity, saturation, types, diagrams, and elementary extensions are
Volume VIII topics. A theory may be called satisfiable in Volume I when it has
at least one model.
\end{remark*}

\begin{dependencies}
\begin{itemize}
  \item \hyperref[def:truth-structure]{Truth in a Structure}
\end{itemize}
\end{dependencies}

\begin{definitionbox}{Definition (Logical Consequence)}
\begin{definition}[Logical Consequence]
\label{def:log-consequence}
Let \(T\) be a set of sentences and let \(\varphi\) be a sentence in the same
language. The sentence \(\varphi\) is a \emph{logical consequence} of \(T\),
written \(T\models\varphi\), exactly when every model of \(T\) is a model of
\(\varphi\).
\end{definition}
\end{definitionbox}

\begin{remark*}[Standard quantified statement]
\[
T\models\varphi
\Longleftrightarrow
\forall \mathcal M\;(\mathcal M\models T \Rightarrow \mathcal M\models\varphi).
\]
\end{remark*}

\begin{remark*}[Predicate reading]
\(\operatorname{LogicalConsequence}(T,\varphi)\) means that every model of
\(T\) is a model of \(\varphi\).
\end{remark*}

\begin{remark*}[Negated quantified statement]
\[
T\not\models\varphi
\Longleftrightarrow
\exists \mathcal M\;(\mathcal M\models T \land \mathcal M\not\models\varphi).
\]
\end{remark*}

\begin{remark*}[Negation predicate reading]
\(\neg\operatorname{LogicalConsequence}(T,\varphi)\) means that some structure
models \(T\) while failing \(\varphi\).
\end{remark*}

\begin{remark*}[Failure modes]
\begin{description}
  \item[Exposition.] Logical consequence fails by a countermodel: a structure satisfying all
sentences of \(T\) but not satisfying \(\varphi\).
\end{description}
\end{remark*}

\begin{remark*}[Failure modes]
\begin{description}
  \item[Exposition.] \[
\underbrace{\mathcal M\models T}_{\text{premises true}}
\qquad
\underbrace{\mathcal M\not\models\varphi}_{\text{conclusion false}}.
\]
\end{description}
\end{remark*}

\begin{remark*}[Interpretation]
Logical consequence is semantic entailment. The premises in \(T\) leave no
room, across structures, for \(\varphi\) to fail.
\end{remark*}

\begin{dependencies}
\begin{itemize}
  \item \hyperref[def:model-theory-set]{Model of a Set of Sentences}
  \item \hyperref[def:theory]{Theory}
\end{itemize}
\end{dependencies}

\begin{definitionbox}{Definition (Logical Equivalence)}
\begin{definition}[Logical Equivalence]
\label{def:log-equiv}
Two sentences \(\varphi\) and \(\psi\) are \emph{logically equivalent}, written
\(\varphi\equiv\psi\), exactly when each is a logical consequence of the other:
\[
\varphi\models\psi
\quad\text{and}\quad
\psi\models\varphi.
\]
\end{definition}
\end{definitionbox}

\begin{remark*}[Standard quantified statement]
\[
\varphi\equiv\psi
\Longleftrightarrow
\forall \mathcal M\;(\mathcal M\models\varphi \Longleftrightarrow \mathcal M\models\psi).
\]
\end{remark*}

\begin{remark*}[Predicate reading]
\(\operatorname{LogicalEquivalence}(\varphi,\psi)\) means that \(\varphi\) and
\(\psi\) have the same truth value in every structure.
\end{remark*}

\begin{remark*}[Negated quantified statement]
\[
\varphi\not\equiv\psi
\Longleftrightarrow
\exists \mathcal M\;\bigl((\mathcal M\models\varphi \land \mathcal M\not\models\psi)
\lor
(\mathcal M\models\psi \land \mathcal M\not\models\varphi)\bigr).
\]
\end{remark*}

\begin{remark*}[Negation predicate reading]
\(\neg\operatorname{LogicalEquivalence}(\varphi,\psi)\) means that some
structure separates the two sentences.
\end{remark*}

\begin{remark*}[Failure modes]
\begin{description}
  \item[Exposition.] Logical equivalence fails when one sentence is true and the other false in the
same structure.
\end{description}
\end{remark*}

\begin{remark*}[Failure modes]
\begin{description}
  \item[Exposition.] \[
\underbrace{\mathcal M\models\varphi \land \mathcal M\not\models\psi}_{\text{first separation}}
\quad\lor\quad
\underbrace{\mathcal M\models\psi \land \mathcal M\not\models\varphi}_{\text{second separation}}.
\]
\end{description}
\end{remark*}

\begin{remark*}[Interpretation]
Logical equivalence is semantic interchangeability at the sentence level. It
records sameness of truth across all structures in the relevant language.
\end{remark*}

\begin{dependencies}
\begin{itemize}
  \item \hyperref[def:log-consequence]{Logical Consequence}
\end{itemize}
\end{dependencies}
